% ============================================================
% Section 2: Numerical Discretization
% ============================================================

\section{Numerical Discretization}
\pagecheck{3--4}

\instructionplaceholder{Derivation overview}{Derive the discrete formulation for both subsystems at a level that demonstrates method ownership: Finite Element Method for structural dynamics and Finite Volume Method for current loading. Keep this section focused on mathematical discretization, not software implementation.}

\subsection{Finite Element formulation for the structure}
\instructionplaceholder{Required evidence}{Describe discretisation of structural domains, interpolation functions, weak form, element matrices, boundary conditions, and external force mapping. Justify element types and approximation spaces for the expected dynamic range.}

\subsection{Finite Volume formulation for current loading}
\instructionplaceholder{Required evidence}{Define control-volume discretisation, flux treatment, pressure-velocity coupling, boundary conditions, and force extraction procedure on the cylinder. Include how pressure and viscous contributions are computed.}

\subsection{Coupling strategy and assumptions}
\instructionplaceholder{Required evidence}{Explain how current-force outputs inform structural loading. Clarify one-way or two-way coupling, data transfer, update frequency, and consistency assumptions. Discuss expected implications on accuracy and computational cost.}

\subsection{Discretization justification and alternatives}
\instructionplaceholder{Required evidence}{For each major discretization choice, state why it was selected, one credible alternative, and the expected trade-off in accuracy, robustness, or computational cost.}
